\documentclass[twocolumn,letterpaper]{article}

\usepackage[T1]{fontenc}
\usepackage{newtxtext}
\usepackage[protrusion=false,expansion=false]{microtype}
\DisableLigatures[f]{encoding=*} 
\usepackage{graphicx}
\usepackage[table]{xcolor}
\usepackage{array}
\usepackage{enumitem}
\usepackage[letterpaper,
            left=46.8pt,right=46.8pt,
            top=46.8pt,textheight=712pt,
            footskip=12.1pt,
            columnsep=18pt]{geometry}
\input{glyphtounicode}
\pdfgentounicode=1

\definecolor{figborder}{RGB}{214,214,214}
\definecolor{tbldark}{RGB}{51,51,53}
\definecolor{tblgray}{RGB}{190,190,192}
\definecolor{tblhead}{RGB}{244,243,248}
\definecolor{tblhi}{RGB}{249,250,252}

% ---------------------------------------------------------------------
% Fonts / sizes
% ---------------------------------------------------------------------
\renewcommand\normalsize{\fontsize{9.55}{10.8}\selectfont}
\newcommand\abssize{\fontsize{8.6}{10.15}\selectfont}
\newcommand\capsize{\fontsize{8.0}{9.4}\selectfont}
\newcommand\refsize{\fontsize{8.1}{9.4}\selectfont}
\newcommand\tblsize{\fontsize{8.1}{9.4}\selectfont}
\newcommand\tblcapsize{\fontsize{8.5}{9.2}\selectfont}
\newcommand\headsize{\fontsize{10}{12}\selectfont}
\newcommand\subheadsize{\fontsize{10.1}{12}\selectfont}

% ---------------------------------------------------------------------
% Paragraph / line-breaking set-up 
% ---------------------------------------------------------------------
\setlength\parindent{10.8pt}
\setlength\parskip{2.3pt}
\setlength\topskip{10.8pt}
\frenchspacing
\hyphenpenalty=10000 \exhyphenpenalty=10000
\pretolerance=-1 \tolerance=10000 \emergencystretch=0pt
\hbadness=10000 \vbadness=10000 \hfuzz=40pt \vfuzz=40pt
\setlength\rightskip{0pt minus 6pt}  % lets a slightly long line shrink instead of re-breaking
\raggedbottom
\setlength\lineskiplimit{-10pt}

\newcommand\lb{\linebreak[4]\relax}                         % justified forced line break
\newcommand\colbreak{{\parfillskip=0pt\par}\newpage\noindent}  % end column mid-paragraph
\newcommand\colend{{\parfillskip=0pt\par}\newpage}              % same, inside a list
\newcommand\textq{\textquotesingle}               % straight apostrophe
\newcommand\tld{\textasciitilde}                    % 

% page number: centred, small
\makeatletter
\renewcommand\@oddfoot{\hfil{\fontsize{8.3}{9}\selectfont\thepage}\hfil}
\renewcommand\@evenfoot{\@oddfoot}
\renewcommand\@oddhead{}\renewcommand\@evenhead{}
\makeatother

% ---------------------------------------------------------------------
% Headings
% ---------------------------------------------------------------------
\newcommand\Section[1]{%
  \par\vspace{5.4pt}\refstepcounter{section}%
  {\headsize\noindent\makebox[\linewidth][c]{\Roman{section}. #1}\par}%
  \vspace{1.25pt}}
\newcommand\SectionNoNum[1]{%
  \par\vspace{5.4pt}%
  {\headsize\noindent\makebox[\linewidth][c]{#1}\par}%
  \vspace{1.25pt}}
\newcommand\Subsection[1]{%
  \par\vspace{3.5pt}\refstepcounter{subsection}%
  {\subheadsize\itshape\noindent\makebox[\linewidth][l]{\Alph{subsection}. #1}\par}%
  \vspace{-0.5pt}}

% ---------------------------------------------------------------------
% Lists
% ---------------------------------------------------------------------
\setlist{nosep,topsep=0pt,partopsep=0pt,parsep=0pt,itemsep=2pt}
\newlist{bul}{itemize}{1}
\setlist[bul]{label=\textbullet,leftmargin=13.3pt,labelindent=2.6pt,
  labelwidth=10.7pt,labelsep=0pt,align=left,topsep=0pt,itemsep=2pt}
\newlist{num}{enumerate}{1}
\setlist[num]{label=\arabic*),leftmargin=13.3pt,labelindent=2.4pt,
  labelwidth=10.9pt,labelsep=0pt,align=left,topsep=0pt,itemsep=2pt}

% -------------------------------------------------------------
\newcommand\eqline[2]{%
  \par\refstepcounter{equation}%
  {\noindent\hspace*{#1}\rlap{\itshape #2}\hfill(\theequation)\par}\vspace{1.6pt}}

\newcommand\colfig[1]{%
  \par{\noindent\makebox[\linewidth][c]{\setlength\fboxsep{0pt}\setlength\fboxrule{0.4pt}%
   \fcolorbox{figborder}{white}{\includegraphics[width=245.3pt]{#1}}}\par}}
\newcommand\figcap[1]{%
  \par{\capsize\noindent\refstepcounter{figure}\textbf{Fig.\ \thefigure.} #1\par}}

% ---------------------------------------------------------------------
% Tables
% ---------------------------------------------------------------------
\renewcommand\thetable{\Roman{table}}
\newcommand\tblcap[1]{\par{\tblcapsize\centering\refstepcounter{table}%
   {\fontsize{8.2}{9.1}\selectfont\textbf{TABLE \thetable}}\\[0.4pt] #1\par}}
\newcolumntype{L}[1]{>{\hspace{1.8pt}}p{\dimexpr#1-1.8pt\relax}}
\newcolumntype{C}[1]{>{\centering\arraybackslash}p{#1}}
\newcolumntype{V}{!{\color{tblgray}\vrule width 0.4pt}}  % light-grey vertical rule
\newcommand\hstrut{\rule[-3.8pt]{0pt}{11.8pt}}             % header row height
\newcommand\rstrut{\rule[-3.32pt]{0pt}{11.12pt}}           % body row height
\newcommand\ghline{\arrayrulecolor{tblgray}\hline}         % light-grey row rule
\newcommand\dhline{\noalign{\begingroup\color{tbldark}\hrule height 0.9pt\endgroup}}  % top/bottom
\newcommand\dhlineh{\noalign{\begingroup\color{tbldark}\hrule height 0.6pt\endgroup}} % under header
\newenvironment{tframe}{%
  \par\noindent\setlength\tabcolsep{0pt}\setlength\arrayrulewidth{0.4pt}%
  \tblsize\begin{lrbox}{0}}{\end{lrbox}\makebox[\linewidth][c]{\usebox{0}}\par}
%

% =====================================================================
\begin{document}
\normalsize

\twocolumn[{%
\centering
\vspace*{7.4pt}
{\fontsize{21.7}{26}\selectfont
OptiPulse: A Noise-Aware Dual-Stream 3D-CNN\\
Framework with Linear-First Region Combining and\\
Causal Spectral Tracking for Remote Heart Rate Monitoring\par}
\vspace{9.2pt}
{\fontsize{10.6}{13.7}\selectfont
Jayant Mitawa (230504), Priyen Parmar (230735), Satwik Raj Wadhwa (230937),\\
Divyansh Bansal (230381), and Anuj Agrawal (230169)\par}
\vspace{3.2pt}
{\fontsize{10.1}{12.5}\selectfont\itshape
Undergraduate Project (UGP) --- Department of Design\\
Under the Supervision of Prof. Amar Kumar Behera\par}
{\fontsize{9.55}{12.4}\selectfont
Indian Institute of Technology Kanpur (IIT Kanpur)\par}
\vspace{10.7pt}
}]

% ============================ PAGE 1, LEFT ============================
{\abssize\bfseries\setlength\parindent{9.5pt}\setlength\parskip{2pt}
\textit{Abstract---} This project presents OptiPulse, a contactless\lb
vital-sign monitoring system for estimating heart rate (HR) in real\lb
time from facial video data. Traditional monitoring approaches\lb
rely on contact-based sensors such as ECG leads and finger pulse\lb
oximeters, which require physical attachment and limit scalability\lb
for continuous remote sensing. While remote\newline
photoplethysmography (rPPG) enables non-contact estimation\lb
from standard RGB cameras, real-world video streams introduce\lb
severe degradations: low-light photon shot and read noise,\lb
MPEG/H.264 compression artifacts, head movement, frame-rate\lb
jitter, and automatic camera exposure adjustments.

To overcome the limitations of our initial PhysFormerX\lb
baseline prototype (which achieved an end-to-end MAE of 6.81\lb
BPM due to hard YCbCr mask flicker, assumed constant frame\lb
rates, and resting-rate bias), we propose a unified six-stage\lb
pipeline shared identically across preprocessing, training, and live\lb
inference. The system integrates asynchronous YOLOv8 face\lb
localization, motion-adaptive bounding-box stabilization,\lb
true-timestamp 30 Hz resampling, and a 10-model heterogeneous\lb
ensemble of two complementary 3D-CNN architectures:\lb
VitalNet-Pix (v1, dual-stream pixel path with soft appearance\lb
attention) and VitalNet-MRC (v4, augmenting the pixel path with\lb
a linear-first float32 multi-scale region branch, time-varying\lb
softmax maximal-ratio combining, and input-noise-gated fusion).

Trained on audited UBFC-rPPG videos with cross-illumination\lb
self-consistency and uniform HR time-stretch resampling (45--180\lb
BPM), the ensemble spectra are decoded by a prior-free causal\lb
Hidden Markov Model (HMM) tracker and filtered by a 5-second\lb
median SNR honesty gate (>= -1.5 dB). Under strict subject-wise\lb
5-fold cross-validation on 43 unseen subjects, OptiPulse achieves\lb
an MAE of 1.45 BPM in normal lighting (96.3\% within 5 BPM, r\lb
= 0.948), while reducing low-light MAE by 30.8\% (10.91 to 7.55\lb
BPM) and motion MAE by 23.2\% (4.06 to 3.12 BPM). In\lb
end-to-end live evaluation across all raw videos, OptiPulse\lb
achieves 1.59 BPM MAE overall and 1.17 BPM MAE on\lb
honesty-gated displayed readings (95\% coverage).

\vspace{2.3pt}
\textit{Index Terms---} remote photoplethysmography (rPPG),\lb
contactless heart rate monitoring, 3D convolutional neural\lb
networks, maximal-ratio combining (MRC), noise-gated fusion,\lb
cross-illumination consistency, Hidden Markov Model (HMM),\lb
YOLOv8 face tracking.\par}

\vspace{8.0pt}
\Section{INTRODUCTION}
Monitoring vital signs such as heart rate (HR) in beats per\lb
minute (BPM) is essential for non-invasive physiological sensing,\lb
human-computer interaction, and remote telemetry. Conventional\lb
monitoring systems rely on contact-based sensors, including\lb
electrocardiography (ECG) devices and optical finger pulse\lb
oximeters, which require direct physical attachment. While\lb
accurate in stationary setups, contact sensors restrict movement
\colbreak
% ============================ PAGE 1, RIGHT ===========================
and limit continuous long-term monitoring on everyday\lb
consumer devices.

Recent advancements in computer vision have enabled\lb
contactless physiological measurement using video-based remote\lb
photoplethysmography (rPPG) [1], [8], [9]. Subtle periodic\lb
variations in facial optical reflectance caused by pulsatile blood\lb
volume changes can be captured using standard RGB cameras\lb
and analyzed over time. With every cardiac cycle, microvascular\lb
blood volume increases transiently; because hemoglobin absorbs\lb
green and blue wavelengths strongly, reflected pixel intensity\lb
dips by approximately 0.2\% to 1.0\% per beat (measured at\lb
0.34\% in our live recordings). Because a single 8-bit quantization\lb
step equals 1/255 (0.39\%) of mid-gray, the pulsatile signal in any\lb
single pixel is smaller than camera quantization and sensor read\lb
noise. Recovering an accurate waveform therefore requires\lb
averaging thousands of facial pixels spatially and integrating\lb
over 10-second windows temporally.

To address the real-world fragility of video-based rPPG, this\lb
Undergraduate Project (UGP)---conducted under the supervision\lb
of Prof. Amar Kumar Behera in the Department of Design at IIT\lb
Kanpur---proposes \textbf{OptiPulse} , a structured end-to-end pipeline\lb
consisting of:
\begin{bul}
\item True-timestamp video acquisition (30--60 FPS) with\lb
duplicate-frame suppression and real-time scene diagnostics.
\item Asynchronous YOLOv8 ROI detection coupled with a\lb
motion-adaptive exponential moving average smoother and a\lb
1.6x enlarged crop.
\item Uniform 30.0 Hz temporal resampling and in-network\lb
6-channel normalization (DiffNormalized motion +\lb
standardized appearance).
\item A 10-model heterogeneous 3D-CNN ensemble combining\lb
\textbf{VitalNet-Pix (v1)} and \textbf{VitalNet-MRC (v4)} trained with\lb
cross-illumination consistency and HR time-stretch\lb
augmentation.
\item Spectral-domain ensemble fusion decoded by a prior-free\lb
causal HMM tracker, slew-rate limiter, and 5-s median SNR\lb
honesty gate.
\end{bul}

\Section{PROBLEM STATEMENT}
The primary problem addressed in this project is the accurate,\lb
unbiased extraction of heart rate across the full 39--192 BPM\lb
band (0.65--3.2 Hz) from standard RGB video streams without\lb
physical sensors. Detailed diagnosis of our initial PhysFormerX\lb
baseline (which had an end-to-end MAE of 6.81 BPM) revealed\lb
five critical real-world challenges:
\newpage
% ============================ PAGE 2, LEFT ============================
\Subsection{Key Challenges}
\begin{num}
\item \textbf{Timestamp and Frame-Rate Drift:} Consumer webcams\lb
rarely deliver an exact 30.0 FPS. Our Apple M4 FaceTime\lb
sessions delivered 24.0 FPS with 137 bursty gaps of\lb
0.10--0.35 s, while DroidCam virtual phone cameras repeated\lb
54\% of frames to pad 28 FPS streams to 60 FPS. Assuming a\lb
nominal 30 FPS scales every FFT frequency bin by 30/24 =\lb
+25\%.
\item \textbf{Hard Mask Flicker \& Train--Live Mismatch:} The earlier\lb
pipeline used a binary YCbCr chrominance threshold whose\lb
boundary pixels flickered in and out of the spatial mean from\lb
frame to frame, injecting step-change noise larger than the\lb
pulse itself.
\item \textbf{Non-Linear Signal Destruction in Low Light:} In dark\lb
scenes where per-pixel pulse SNR falls below -40 dB,\lb
applying BatchNorm and ELU non-linearities at full pixel\lb
resolution before spatial averaging destroys the weak periodic\lb
signal.
\item \textbf{Dataset HR Distribution Bias (70--100 BPM Collapse):}\lb
Because UBFC-rPPG labels cluster around resting rates,\lb
standard regression and unconstrained networks gravitate\lb
toward \tld75 BPM and fail at elevated heart rates (120--170\lb
BPM).
\item \textbf{Camera Auto-Exposure \& False Tracker Locks:} Backlit\lb
scenes trigger webcam auto-exposure swings of -11\% to +8\%\lb
(\tld30x larger than the 0.34\% pulse). A temporal tracker can\lb
lock onto steady camera noise on a blank wall and report\lb
high self-consistency unless gated by true spectral SNR.
\end{num}

\Subsection{Objective}
The objective of OptiPulse is to develop a real-time\lb
contactless monitoring system across 39--192 BPM that: (i)\lb
achieves \tld1.5 BPM MAE on unseen subjects in normal light; (ii)\lb
reduces error by >25\% under low light and head motion; (iii)\lb
guarantees zero train--live discrepancy via a unified library (\lb
\textbf{optipulse/} ); (iv) runs smoothly on NVIDIA CUDA GPUs, Apple\lb
Silicon MPS GPUs, and CPUs; and (v) suppresses unreliable\lb
readings via an empirically calibrated SNR honesty gate.

\Section{ARCHITECTURAL OVERVIEW}
OptiPulse is organized into a single-source-of-truth Python\lb
package ( \textbf{optipulse/} ) shared identically by offline dataset\lb
preprocessing, 5-fold cross-validation training, deterministic\lb
stress evaluation, and the real-time webcam application (Fig. 1).\lb
Separating asynchronous face detection from 30 Hz resampled\lb
3D-CNN inference ensures high frame throughput and accurate\lb
frequency estimation.
\newpage
% ============================ PAGE 2, RIGHT ===========================
\vspace*{-8.4pt}
\colfig{figures/optipulse_pipeline_diagram.png}
\vspace{-1.5pt}
\figcap{The six-stage OptiPulse modular pipeline integrating true-timestamp\lb
capture, YOLOv8 tracking with motion-adaptive BoxSmoother, uniform 30\lb
Hz resampling, the 10-model VitalNet-Pix and VitalNet-MRC ensemble,\lb
spectral fusion with causal HMM tracking, and the SNR honesty gate.}

\vspace{3.2pt}
\Subsection{Asynchronous YOLOv8 Detection Pipeline}
Each incoming BGR frame is stamped immediately upon\lb
capture with its monotonic clock time via \textbf{time.perf\_counter()} in\lb
a dedicated background capture thread (45-frame queue). Exact\lb
duplicate frames from virtual camera drivers are detected via an\lb
8x8 subsampled equality check and skipped before buffering.\lb
Facial localization uses our custom-trained YOLOv8 model (\lb
\textbf{yolov8\_face\_nose\_chest.pt} , input size 320x320, confidence\lb
threshold 0.30) selecting the highest-confidence face bounding\lb
box. On CUDA GPUs, detection runs every 2nd frame at 30 FPS\lb
(every 4th at 60 FPS); on macOS MPS and CPU machines, an\lb
\textbf{AsyncFace} background thread runs YOLOv8 at 4--8 Hz so\lb
detector latency (\tld147 ms on CPU) never blocks frame capture.\lb
If YOLOv8 is unavailable, the detector automatically falls back\lb
to OpenCV\textq s frontal Haar cascade.

\Subsection{Motion-Adaptive BoxSmoother \& 1.6x Crop}
Even with an accurate detector, bounding-box coordinates\lb
exhibit 1--3 pixel jitter between frames, shifting borders and\lb
creating false intensity oscillations. Fixed heavy EMA\lb
smoothing, however, lags behind real head movement. We\lb
implement \textbf{BoxSmoother} , a causal motion-adaptive EMA filter\lb
over the square box state s\_t = (c\_x, c\_y, S). Given relative\lb
center displacement d\_t = \textbar\textbar m\_t - s\_\{t-1\}\textbar\textbar\ / S\_\{t-1\}, the\lb
smoothing gain a\_t scales linearly between a\_slow = 0.12 (strong\lb
jitter suppression when stationary) and a\_fast = 0.75 (rapid\lb
tracking when displacement reaches 8\% of face size), while\lb
scale smoothing is capped at a\_S = min(a\_t, 0.25) to prevent\lb
zoom-pumping:
\vspace{3.7pt}
\eqline{28pt}{s\_t = (1 - a\_t) s\_\{t-1\} + a\_t m\_t,\ \ \ a\_t in [0.12, 0.75]}
\newpage
% ============================ PAGE 3, LEFT ============================
\hspace*{10.8pt}Using s\_t, \textbf{crop\_square()} warps a square of side 1.6 \textit{S\_t in a\lb
single sub-pixel affine step ( cv2.warpAffine with replicate\lb
border padding) to an 80x80 RGB crop. Crucially, no hard\lb
YCbCr mask is used} ; eliminating binary chrominance\lb
thresholding removes mask-flicker noise and allows the neural\lb
network to learn a smooth, soft spatial attention mask.

\Subsection{True-Timestamp Bin Averaging \& 30 Hz Resampling}
Every 0.5 s (once >= 6.0 s of crops are buffered, up to a 10.0\lb
s window; any gap > 1.0 s resets the buffer), the buffered\lb
sequence of (timestamp, crop) pairs is converted to a uniform\lb
30.0 Hz clip (300 frames for 10 s). When the camera delivers >\lb
39 FPS (such as 60 FPS phone streams via DroidCam),\lb
\textbf{bin\_average()} averages all frames falling into each 1/30 s slot\lb
before linear timestamp interpolation ( \textbf{resample\_uniform()} ),\lb
reducing per-frame sensor noise by sqrt(k). Finally,\lb
\textbf{to\_model\_input()} applies a deterministic center zoom of 0.93\lb
(80x80 -> 72x72 RGB in float32), matching the center of the\lb
training zoom distribution [0.86, 1.00].

\Section{VITALNET ARCHITECTURES \& LEARNING}
OptiPulse ensembles two complementary 3D convolutional\lb
neural network architectures implemented in \textbf{optipulse/model.py}\lb
(Fig. 2): \textbf{VitalNet-Pix (v1)} and \textbf{VitalNet-MRC (v4)} .

\vspace{3.7pt}
\colfig{figures/vitalnet_architecture_diagram.png}
\vspace{-0.8pt}
\figcap{Architecture of VitalNet-Pix (v1) and VitalNet-MRC (v4). Both\lb
share in-network 6-channel normalization and dual 3D-CNN towers with\lb
soft spatial attention (g1, g2). VitalNet-MRC adds a linear-first float32\lb
multi-scale RegionBranch (46 regions across 6x6, 3x3, 1x1 grids),\lb
time-varying softmax maximal-ratio combining, and an\newline
input-noise-conditioned MLP fusion gate.}

\vspace{2.8pt}
\Subsection{In-Network 6-Channel Difference-Frame Representation}
To ensure training and live inference can never normalize\lb
inputs differently, \textbf{normalize\_inputs()} executes inside\lb
\textbf{VitalNet.forward()} . Given RGB clip X in [0, 1] of shape [B, 3,\lb
T, 72, 72], it constructs a 6-channel tensor Z = [D \textbar\textbar\ A] of shape\lb
[B, 6, T, 72, 72] combining a 3-channel \textbf{DiffNormalized motion\lb
stream} D and a 3-channel \textbf{standardized appearance stream} A\lb
[4], [14]:
\vspace{2.3pt}
\eqline{9.2pt}{D\_t = clip( ((X\_t - X\_\{t-1\}) / (X\_t + X\_\{t-1\} + eps)) / std, -6, 6 )}

where eps = 2/255 and A = (X - mean(X)) / (std(X) + 1e-4).\lb
Because D\_t is a relative frame-to-frame intensity ratio, global\lb
brightness scaling (X -> kX) cancels out analytically, providing\lb
built-in invariance to dim illumination while removing static\lb
appearance.

\Subsection{VitalNet-Pix (Version 1): Dual-Stream 3D-CNN with Soft Attention}
VitalNet-Pix (0.600M parameters, base width c = 32)\lb
processes the motion and appearance streams through parallel\lb
3D-CNN towers (PhysNet/TS-CAN style [4], [5]) that preserve\lb
full temporal length T = 300 throughout:
\begin{bul}
\item \textbf{Spatial Stems (72x72 -> 36x36):} Conv3d(3->16, kernel
\colend
\end{bul}
% ============================ PAGE 3, RIGHT ===========================
\begin{bul}
\item[] 1x5x5, stride 1x2x2) + BatchNorm3d + ELU on both streams.
\item \textbf{Stage 1 (36x36, 32 ch) \& Soft Attention Mask g1:} Two\lb
spatiotemporal ConvBlock3d(3x3x3) layers on the motion\lb
stream and two spatial ConvBlock3d(1x3x3) layers on the\lb
appearance stream. An \textbf{Attention} module applies a 1x1x1\lb
Conv3d and sigmoid normalized to spatial mean 0.5,\lb
multiplying the motion features (M\_1 <- M\_1 g\_1) before\lb
AvgPool3d(1x2x2).
\item \textbf{Stage 2 (18x18 -> 9x9, 64 ch) \& Stage 3 (9x9, 64 ch):} Stage\lb
2 applies dual-stream convolutions, soft attention mask g\_2,\lb
and spatial pooling to 9x9. Stage 3 applies two temporally\lb
dilated 3D convolutions (d\_t = 2, 4), 3D dropout (p = 0.10),\lb
and spatial mean pooling to yield pixel features f\_pix of\lb
shape [B, 64, T].
\item \textbf{Dilated Temporal Residual Head:} Four 1D residual blocks (\lb
\textbf{TemporalRes} , kernel 5, dilations d = 1, 2, 4, 8, dropout\lb
0.20) followed by a 1x1 Conv1d produce the zero-mean\lb
rPPG waveform y of shape [B, T].
\end{bul}

\Subsection{VitalNet-MRC (Version 4): Linear-First Region Branch \& Noise Gate}
Our Fold-0 diagnosis revealed that while darkening alone\lb
(x0.10) raised MAE modestly (1.06 -> 1.59 BPM) and sensor\lb
noise alone gave 2.51 BPM, their combination caused MAE to\lb
jump to 8.78 BPM: when per-pixel pulse sits \tld40 dB below\lb
sensor noise, applying BatchNorm and ELU non-linearities at\lb
36x36 resolution destroys the weak pulse before pixels are\lb
averaged. \textbf{VitalNet-MRC (v4)} (0.618M parameters) solves this\lb
with three architectural innovations inspired by DistancePPG [8]\lb
and RhythmNet [11]:
\begin{num}
\item \textbf{Linear-First Multi-Scale Region Integration in float32:}\lb
Using Stage-1 attention g\_1 as a spatial weight, raw RGB\lb
pixels are averaged \textit{linearly in float32} (disabling bfloat16\lb
autocast, since bfloat16\textq s 0.4\% precision is coarser than the\lb
0.3\% pulse) across three multi-scale grids: 6x6 (36 regions),\lb
3x3 (9 regions), and 1x1 (1 global region), totaling R = 46\lb
regions. Averaging N pixels linearly reduces independent\lb
sensor noise by sqrt(N) prior to any non-linearity.
\item \textbf{Time-Varying Softmax Maximal-Ratio Combining\lb
(MRC):} Each region trace is diff-normalized, whitened by its\lb
own noise scale s\_m, concatenated with its relative log-noise\lb
level, and passed through a shared dilated 1D-CNN (d = 1, 2,\lb
4). A scoring Conv1d + softmax across the 46 regions\lb
computes time-varying reliability weights w\_\{m,t\} to form\lb
the pooled region representation f\_reg:
\end{num}
\vspace{1.9pt}
\eqline{0pt}{w\_\{m,t\} = softmax\_m( score(h\_\{m,t\}) ),\ \ f\_\{reg,t\} = sum\_m w\_\{m,t\} h\_\{m,t\}}
\vspace{0.2pt}
\begin{num}[start=3]
\item \textbf{Input-Noise-Gated Two-Path Fusion:} The network\lb
measures its own input noise via per-pixel 2nd-difference\lb
log-noise n\_pix and mean region log-noise n\_reg. A 2-layer\lb
MLP maps [n\_pix, n\_reg] to gates [g\_pix, g\_reg] = 2\lb
sigmoid(MLP(n\_pix, n\_reg)), weighting f\_pix against f\_reg\lb
via a 1x1 Conv1d + BN + ELU fusion layer before the\lb
temporal residual head.
\end{num}

\Subsection{Training Objective \& Cross-Illumination Self-Consistency}
Both model families optimize a supervised composite loss\lb
combining negative Pearson correlation (waveform shape\lb
fidelity) and frequency-domain Spectral KL divergence over a\lb
fine 1-BPM explicit DFT grid (39--192 BPM):
\vspace{1.7pt}
\eqline{6.1pt}{L\_sup = L\_Pearson(y\_hat, y) + 0.5 * KL( PSD(y) \textbar\textbar\ PSD(y\_hat) )}
\vspace{0.2pt}

For \textbf{VitalNet-MRC} , we introduce \textbf{cross-illumination\lb
self-consistency} (pair\_w = 0.5): each training clip is shown twice\lb
with identical geometry and motion---once well-lit (X\_lit) and
\colbreak
% ============================ PAGE 4, LEFT ============================
once forced into low light with shot/read noise (X\_dark). The\lb
dark prediction must match the stop-gradient waveform\lb
extracted by the model itself from the well-lit view:
\vspace{1.7pt}
\eqline{0pt}{L\_tot = L\_sup(lit) + L\_sup(dark) + 0.5 * L\_Pearson(y\_dark, sg(y\_lit))}
\vspace{0.2pt}

Unlike finger contact-PPG (which has a finger-to-face pulse\lb
transit lag and different dicrotic shape), stopgrad(y\_lit) shares the\lb
exact facial pulse phase and morphology, isolating illumination\lb
and sensor noise.

\Subsection{Unbiased HR Time-Stretch \& Webcam-Realistic GPU Augmentation}
To prevent the network from memorizing the 65--100 BPM\lb
training distribution, \textbf{BatchBuilder} applies \textbf{HR time-stretch\lb
resampling} to 60\% of clips: a target HR is drawn uniformly\lb
from 45--180 BPM and both video and PPG are time-stretched\lb
by factor r in [0.55, 1.80]. On-GPU augmentations (\lb
\textbf{optipulse/augment.py} ) further simulate low light (gain down to\lb
6\%, gamma, shot/read noise, 8-bit quantization), color cast, side\lb
light, auto-exposure drift/steps, quasi-periodic head motion and\lb
jumps, blur, and 10--25 FPS frame drops re-interpolated to 30\lb
Hz.

\Section{DEPLOYMENT \& REAL-TIME OPTIMIZATION}
\vspace{1.1pt}
\Subsection{Spectral Ensemble Fusion \& Harmonic SNR}
At inference, the 10 fold checkpoints (5 \textbf{pix\_fold0..4} + 5\lb
\textbf{mrc\_fold0..4} ) run in ensemble. Because models can predict\lb
waves with opposite sign or slight phase shifts, averaging raw\lb
waveforms causes cancellation. Instead, each waveform is\lb
detrended (Tarvainen smoothness priors, lambda = 100 [10]),\lb
bandpass-filtered (2nd-order Butterworth, 0.65--3.2 Hz = 39--192\lb
BPM), and transformed via a 4096-point Hann FFT. Spectral\lb
SNR (in dB) is computed as the ratio of power within +-0.1 Hz\lb
of the peak frequency f\_0 and +-0.2 Hz of its first harmonic\lb
2f\_0 to all remaining in-band power [9].

\Subsection{Prior-Free Causal HMM Heart-Rate Tracker}
In severe low light, a random noise peak beats the true\lb
cardiac peak in \tld35\% of single 10-s windows. However, the true\lb
heart rate is continuous over time while noise peaks jump\lb
randomly. \textbf{HRTracker} ( \textbf{optipulse/tracking.py} ) decodes the\lb
ensemble\textq s band-normalized PSD on a 1-BPM grid (39--192\lb
BPM) using a causal forward HMM filter: (i) a \textbf{strictly uniform\lb
initial prior} (no pull toward 70--80 BPM); (ii) a Gaussian\lb
random-walk transition kernel (sigma = 4.0 BPM/s) that follows\lb
exercise recovery while rejecting sudden jumps; and (iii)\lb
\textbf{overlap-tempered spectral evidence} L\_t(b) = ((1 - eps)\lb
PSD(b)\textasciicircum 2 + eps/154)\textasciicircum (dt / 1.0s), which down-weights 0.5-s\lb
overlapping window updates so evidence is never\lb
double-counted. Tracker confidence C\_t in [0, 1] is the posterior\lb
probability mass within +-3 BPM of the MAP peak.

\Subsection{Display Lock, Slew Limit, Honesty Gate \& Cross-Platform UI}
To ensure reliable live operation, \textbf{HREngine} applies three\lb
display rules: (1) \textbf{Confidence Lock:} no reading is displayed until\lb
C\_t >= 0.60 (typically 6.5--12 s); (2) \textbf{Physiological Slew Limit:}\lb
displayed BPM changes by at most 3.0 BPM/s, or 9.0 BPM/s\lb
when C\_t >= 0.75; and (3) \textbf{5-s Median SNR Honesty Gate:} if\lb
the 5-s median SNR falls below -1.5 dB, the readout is hidden (\lb
\textit{\textq Pulse signal too weak\textq} ), lowering displayed error on UBFC\lb
from 1.59 to 1.17 BPM while hiding only 5\% of updates. As\lb
shown in Fig. 3, \textbf{LiveUI} renders the camera feed with\lb
state-colored brackets, a fast 96x96 attention glow, live BPM,\lb
beating heart, 5-bar SNR indicator, 60-s HR trend, pulse\lb
waveform, and \textbf{SceneMonitor} backlighting/motion alerts. On
\colbreak
% ============================ PAGE 4, RIGHT ===========================
macOS ( \textbf{launchers/macos/run\_mac.sh} ) and Windows (\lb
\textbf{launchers/windows/run\_optipulse.bat} ), \textbf{pick\_device()}\lb
automatically selects CUDA or verifies Apple Silicon MPS\lb
against CPU within 2\% tolerance.

\vspace{0.5pt}
\colfig{figures/live_ui_preview.png}
\vspace{-1.5pt}
\figcap{The OptiPulse real-time instrument panel (optipulse/ui.py) showing\lb
face tracking brackets, the learned 36x36 attention heatmap glow on the\lb
forehead and cheeks, live BPM readout, state chip, 5-bar SNR indicator,\lb
60-s HR trend, and filtered rPPG waveform.}

\vspace{3.7pt}
\Section{DATASET PREPARATION \& GROUND-TRUTH AUDIT}
We evaluate on all 50 videos of the public \textbf{UBFC-rPPG}\lb
dataset (Dataset 1: 8 videos; Dataset 2: 42 subjects, 640x480\lb
MPEG-4 at \tld0.8 Mbit/s with finger pulse-oximeter PPG ground\lb
truth). Inspecting Dataset 2\textq s raw ground-truth timestamps\lb
revealed an 11 ms / 44 ms logging-timer quantization jitter;\lb
\textbf{optipulse/ubfc.py} removes this jitter via linear regression of\lb
frame index against logged time, recovering true capture rates\lb
(23.9--30.1 FPS). Running \textbf{pipeline/preprocess\_ubfc.py} extracts\lb
1.6x smoothed YOLOv8 face crops (80x80 RGB) and resamples\lb
both video and contact PPG to an exact 30.0 Hz timeline (\lb
\textbf{100.0\% face detection rate} across all 50 videos).

Before training, \textbf{pipeline/audit\_dataset.py} audited every\lb
video\textq s ground truth across 10-s windows (Fig. 4). Seven subjects\lb
failed the quality audit (oximeter-vs-PPG HR disagreement > 10\lb
BPM in > 30\% of windows, or median PPG SNR < 2.0 dB):\lb
\textbf{sub11, sub18, sub20, sub24} (oximeter HR dropouts with usable\lb
PPG waveforms) and \textbf{sub3, sub4, sub14} (noisy contact PPG). All\lb
7 flagged subjects are excluded from training across all folds so\lb
corrupted labels never degrade the network, and main CV\lb
metrics are reported on the \textbf{43 reliable subjects} (with the 7\lb
held-out subjects evaluated separately). Auditing the legacy\lb
\textbf{preprocessed\_vidformer2} cache also revealed that old windows\lb
had been saved with a wrong hard-coded 28.0 FPS instead of\lb
28.67 FPS.
\newpage
% ============================ PAGE 5, LEFT ============================
\vspace*{-8.4pt}
\colfig{figures/audit_hr_timelines.png}
\vspace{-0.5pt}
\figcap{Automated ground-truth quality audit across all 50 UBFC-rPPG\lb
subjects (pipeline/audit\_dataset.py), comparing contact-PPG spectral HR\lb
(black), oximeter sensor HR (cyan dashed), and video POS HR (red dotted).\lb
Seven flagged subjects are excluded from training.}

\vspace{3.6pt}
\Section{EXPERIMENTAL DEPTH AND TRIED METHODS}
Before finalizing the OptiPulse architecture, we systematically\lb
implemented and benchmarked classical non-learning methods,\lb
deep spatiotemporal baselines from our initial exploration phase,\lb
and four targeted architectural ablations on Fold 0 (Table I):
\begin{bul}
\item \textbf{Classical Non-Learning Baselines (GREEN, CHROM,\lb
POS):} Spatial RGB means inside a central facial ellipse\lb
projected via GREEN, CHROM [9], and POS [6] yield 15.53,\lb
17.36, and 20.56 BPM MAE on UBFC because 0.8 Mbit/s\lb
MPEG-4 compression corrupts fixed color ratios.
\item \textbf{Earlier Deep Architectures (DeepPhys, PhysNet,\lb
HybridRPPG, TempResNet, AdvancedPhysNet,\lb
PhysFormerX):} In our Phase-1 exploration on UBFC-rPPG,\lb
raw DeepPhys (26.45 BPM), PhysNet (44.99 BPM),\lb
HybridRPPG (18.64 BPM), TempResNet (12.71 BPM),\lb
Lightweight EMD-CNN (12.02 BPM), VidFormer (9.63\lb
BPM), AdvancedPhysNet (8.23 BPM), and PhysFormerX\lb
(6.28 BPM window-split / 6.81 BPM end-to-end live) were\lb
limited by hard YCbCr mask flicker, timestamp drift, and\lb
lack of HR time-stretch augmentation.
\item \textbf{Fold-0 Development Ablations (Rejected vs. Retained):} On\lb
Fold 0 (10 unseen subjects), increasing low-light\lb
augmentation strength and training for 40 epochs ( \textbf{Run 2} )\lb
degraded motion MAE from 3.76 to 8.25 BPM ( \textit{rejected} ).\lb
Inserting a [1,2,1]/4 temporal pre-filter stem before frame\lb
differencing ( \textbf{VitalNet v3} ) hurt normal MAE (1.06 -> 2.14\lb
BPM) and low-light MAE (8.61 -> 11.23 BPM) ( \textit{rejected} ).
\colbreak
\end{bul}
% ============================ PAGE 5, RIGHT ===========================
\begin{bul}
\item[] Test-time Gaussian spatial blur degraded low-light MAE to\lb
14--18 BPM ( \textit{rejected} ). In contrast, combining \textbf{VitalNet-Pix\lb
(v1)} + \textbf{VitalNet-MRC (v4)} with spectral fusion and causal\lb
HMM tracking lowered Fold-0 clean MAE to \textbf{1.01 BPM} ,\lb
motion MAE to \textbf{2.37 BPM} , and low-light + motion MAE to\lb
\textbf{5.92 BPM} ( \textit{retained} ).
\item \textbf{Real-World Blank-Wall Control Experiment (Apple M4\lb
Mac):} During live testing on an Apple M4 MacBook in\lb
backlit conditions (face brightness 80 vs. background 163,\lb
auto-exposure swings of -11\% to +8\%, SNR -2.1 dB),\lb
running the exact same 10-model ensemble on a \textbf{blank wall\lb
and T-shirt patch} produced an identical SNR (-2.4 dB) and\lb
high tracker confidence (0.80). This control experiment\lb
proved that when signal SNR < -1.5 dB, a single 10-s\lb
window cannot be distinguished from camera auto-exposure\lb
noise---motivating our -1.5 dB SNR honesty gate and\lb
real-time \textbf{SceneMonitor} backlighting warnings.
\end{bul}

\vspace{4.6pt}
\tblcap{FOLD-0 DEVELOPMENT ABLATION (10 UNSEEN SUBJECTS, MAE\\
IN BPM) \& PHASE-1 BASELINE COMPARISON}
\vspace{0.6pt}
\begin{tframe}
\begin{tabular}{@{}L{119.0pt}V C{31.86pt}V C{31.86pt}V C{31.86pt}V C{31.5pt}@{}}
\rowcolor{tblhead}\textbf{Model / Ablation Variant}\hstrut & \textbf{Clean} & \textbf{LowLight} & \textbf{Motion} & \textbf{LL+Mot}\\
\dhline
Phase-1: DeepPhys / PhysNet\rstrut & 26.45 & --- & --- & ---\\ \ghline
Phase-1: HybridRPPG / TempResNet\rstrut & 18.64 & --- & --- & ---\\ \ghline
Phase-1: AdvancedPhysNet + YOLO\rstrut & 8.23 & --- & --- & ---\\ \ghline
Phase-1: Legacy PhysFormerX (Live)\rstrut & 6.81 & --- & --- & ---\\ \ghline
Fold-0: Run 1 (VitalNet-Pix v1)\rstrut & 1.06 & 8.61 & 3.76 & 9.77\\ \ghline
Fold-0: Run 2 (Harder Aug, Rejected)\rstrut & 1.08 & 9.00 & 8.25 & 13.00\\ \ghline
Fold-0: VitalNet v3 (Pre-Filter, Rej.)\rstrut & 2.14 & 11.23 & 3.91 & 13.52\\ \ghline
Fold-0: VitalNet v4 (Region-MRC)\rstrut & 1.86 & 7.29 & 3.57 & 9.24\\ \ghline
Fold-0: VitalNet-MRC (v4 + Pair Loss)\rstrut & 1.98 & 7.92 & 4.14 & 11.44\\ \ghline
\rowcolor{tblhi}\textbf{Fold-0: Pix + MRC + HMM Tracked}\rstrut & \textbf{1.01} & \textbf{6.31} & \textbf{2.37} & \textbf{5.92}\\
\end{tabular}
\end{tframe}

\vspace{1.4pt}
\Section{RESULTS}
We evaluated OptiPulse using strict \textbf{subject-wise 5-fold\lb
cross-validation} on UBFC-rPPG (50 videos; fixed 30-epoch\lb
schedule, last epoch kept, zero checkpoint picking on test data).\lb
Every subject is scored exclusively by models that never saw\lb
that person during training. Table II and Table III summarize the\lb
complete 5-fold cross-validation results across the 43 unseen\lb
subjects with reliable ground truth (270 non-overlapping 10-s\lb
windows per condition), illustrated in Fig. 5, Fig. 6, and Fig. 7.

\vspace{1.0pt}
\tblcap{MAIN 5-FOLD SUBJECT-WISE CROSS-VALIDATION MAE (BPM) ON\\
UNSEEN SUBJECTS (N = 43, 270 WINDOWS)}
\vspace{0.6pt}
\begin{tframe}
\begin{tabular}{@{}L{103.8pt}V C{26.82pt}V C{29.34pt}V C{26.82pt}V C{29.34pt}V C{29.7pt}@{}}
\rowcolor{tblhead}\textbf{Method (5-Fold Unseen CV)}\hstrut & \textbf{Norm} & \textbf{LowL} & \textbf{Mot} & \textbf{LL+M} & \textbf{12fps}\\
\dhline
GREEN (Classical)\rstrut & 15.53 & 21.94 & 36.21 & 30.41 & 17.14\\ \ghline
CHROM (Classical [9])\rstrut & 17.36 & 22.45 & 23.51 & 21.96 & 18.60\\ \ghline
POS (Classical [6])\rstrut & 20.56 & 25.97 & 29.17 & 27.64 & 22.74\\ \ghline
Legacy PhysFormerX Monitor\rstrut & 6.81 & --- & --- & --- & ---\\ \ghline
VitalNet-Pix (v1, Single)\rstrut & \textbf{1.42} & 10.91 & 4.06 & 11.47 & 1.50\\ \ghline
VitalNet-MRC (v4, Single)\rstrut & 1.89 & 10.09 & 4.42 & 11.20 & 1.85\\ \ghline
VitalNet-Pix + HMM Tracking\rstrut & 1.56 & 8.02 & 3.48 & 9.92 & 1.55\\ \ghline
VitalNet-MRC + HMM Tracking\rstrut & 1.81 & 7.32 & 3.99 & 9.02 & 1.66\\ \ghline
Pix + MRC Spectral Fusion\rstrut & 1.46 & 9.77 & 3.59 & 10.72 & \textbf{1.43}\\ \ghline
\rowcolor{tblhi}\textbf{OptiPulse Full System}\rstrut & \textbf{1.45} & \textbf{7.55} & \textbf{3.12} & \textbf{8.23} & \textbf{1.59}\\
\end{tabular}
\end{tframe}
\newpage
% ============================ PAGE 6, LEFT ============================
\tblcap{FULL-SYSTEM DETAILED 5-FOLD CV METRICS ACROSS STRESS\\
CONDITIONS (N = 43 UNSEEN SUBJECTS)}
\vspace{0.6pt}
\begin{tframe}
\begin{tabular}{@{}L{88.65pt}V C{29.34pt}V C{31.86pt}V C{31.86pt}V C{29.34pt}V C{34.75pt}@{}}
\rowcolor{tblhead}\textbf{Condition}\hstrut & \textbf{MAE} & \textbf{RMSE} & \textbf{MAPE} & \textbf{r} & \textbf{<=5 BPM}\\
\dhline
\rowcolor{tblhi}\textbf{Normal Light}\rstrut & \textbf{1.45} & \textbf{5.28} & \textbf{2.0\%} & \textbf{0.948} & \textbf{96.3\%}\\ \ghline
Low Light (x0.1 + Noise)\rstrut & 7.55 & 13.89 & 9.1\% & 0.629 & 66.3\%\\ \ghline
Head Motion (Sway/Jumps)\rstrut & 3.12 & 8.14 & 3.9\% & 0.871 & 87.8\%\\ \ghline
Low Light + Motion\rstrut & 8.23 & 14.67 & 9.4\% & 0.557 & 64.8\%\\ \ghline
12 FPS Webcam\rstrut & 1.59 & 5.15 & 2.1\% & 0.951 & 95.2\%\\
\end{tabular}
\end{tframe}

\vspace{1.0pt}
\colfig{figures/cv_scatter_bland_altman.png}
\vspace{-1.5pt}
\figcap{5-fold CV scatter plot and Bland--Altman agreement on unseen\lb
reliable-GT subjects (n = 43, 270 windows, MAE = 1.45 BPM, Pearson r =\lb
0.948, 96.3\% within 5 BPM).}

\vspace{0.5pt}
\colfig{figures/robustness_mae.png}
\vspace{-1.5pt}
\figcap{Stress-test robustness comparison (MAE in BPM, lower is better)\lb
across all five deterministic conditions on unseen subjects (n = 43).}

\vspace{0.5pt}
\colfig{figures/training_curves.png}
\vspace{-1.5pt}
\figcap{Held-out test-fold MAE monitored across the fixed 30-epoch training\lb
schedule for all 5 folds of VitalNet-Pix and VitalNet-MRC.}
\newpage
% ============================ PAGE 6, RIGHT ===========================
\tblcap{END-TO-END LIVE APP EVALUATION ON RAW VIDEOS (0.5S\\
UPDATES) \& SPECTRAL SNR CALIBRATION}
\vspace{0.6pt}
\begin{tframe}
\begin{tabular}{@{}L{124.2pt}V C{39.42pt}V C{41.58pt}V C{41.94pt}@{}}
\rowcolor{tblhead}\textbf{Live Pipeline / SNR Regime}\hstrut & \textbf{Updates} & \textbf{MAE (BPM)} & \textbf{<=5 BPM}\\
\dhline
Legacy PhysFormerX (All 50 Videos)\rstrut & 5,718 & 7.71 & 60.5\%\\ \ghline
Legacy PhysFormerX (43 Reliable)\rstrut & 4,969 & 6.81 & 63.4\%\\ \ghline
\rowcolor{tblhi}\textbf{OptiPulse Live (All 50 Videos)}\rstrut & \textbf{6,055} & \textbf{1.72} & \textbf{94.6\%}\\ \ghline
\rowcolor{tblhi}\textbf{OptiPulse Live (43 Reliable)}\rstrut & \textbf{5,268} & \textbf{1.59} & \textbf{95.9\%}\\ \ghline
\rowcolor{tblhi}\textbf{OptiPulse + Gate (95\% Shown)}\rstrut & \textbf{5,005} & \textbf{1.17} & \textbf{97.4\%}\\ \ghline
SNR < -3 dB (Noise-Dominated)\rstrut & 223 win & 14.30 & 43.5\%\\ \ghline
SNR -3 to 0 dB (Borderline)\rstrut & 429 win & 7.76 & 69.7\%\\ \ghline
SNR 0 to +3 dB (Good Signal)\rstrut & 315 win & \textbf{1.75} & \textbf{95.9\%}\\ \ghline
SNR > +3 dB (High-Quality Pulse)\rstrut & 383 win & \textbf{0.56} & \textbf{99.8\%}\\
\end{tabular}
\end{tframe}

\textbf{Key Experimental Findings:} (1) In normal light, OptiPulse\lb
achieves \textbf{1.45 BPM MAE} (96.3\% of windows within 5 BPM, r =\lb
0.948), reaching the +-1 BPM physical accuracy floor of the\lb
contact finger-PPG reference. (2) Under stress conditions,\lb
combining VitalNet-Pix + VitalNet-MRC with causal HMM\lb
tracking cuts low-light MAE by \textbf{30.8\%} (10.91 -> 7.55 BPM) and\lb
motion MAE by \textbf{23.2\%} (4.06 -> 3.12 BPM). (3) On the 4\lb
never-trained subjects with verified PPG waveforms ( \textbf{sub11,\lb
sub18, sub20, sub24} ), the 10-model ensemble achieves \textbf{0.79\lb
BPM MAE} in normal light and \textbf{1.16 BPM MAE} under motion.\lb
(4) Running the full live application frame-by-frame across all\lb
raw videos ( \textbf{pipeline/compare\_pipelines.py} , Table IV)\lb
achieves \textbf{1.59 BPM MAE} (median subject 1.02 BPM), and\lb
applying the -1.5 dB honesty gate lowers displayed MAE to \textbf{1.17\lb
BPM} across 95\% of updates---a \textbf{4.3x error reduction} over the\lb
legacy PhysFormerX baseline (6.81 BPM).

\Section{CONCLUSION AND FUTURE WORK}
The OptiPulse project demonstrates that achieving\lb
near-1-BPM contactless heart-rate accuracy requires\lb
co-designing signal acquisition, physics-informed neural\lb
architectures, and probabilistic spectral tracking. By integrating\lb
true-timestamp 30 Hz resampling, motion-adaptive YOLOv8\lb
ROI stabilization, in-network 6-channel normalization, a\lb
heterogeneous ensemble of pixel-attention ( \textbf{VitalNet-Pix} ) and\lb
linear-first region-MRC ( \textbf{VitalNet-MRC} ) 3D-CNNs trained\lb
with cross-illumination consistency and uniform HR\lb
time-stretching, prior-free causal HMM tracking, and an SNR\lb
honesty gate, OptiPulse achieves 1.45 BPM CV MAE and 1.17\lb
BPM gated live MAE across unseen subjects.

\Subsection{Key Takeaways}
\begin{bul}
\item \textbf{Zero Train--Live Divergence \& True Timestamps:}\lb
Embedding normalization in \textbf{forward()} and resampling onto\lb
real capture timestamps eliminates frame-rate scaling errors.
\item \textbf{Linear-First Region Integration Beats Pixel Noise:}\lb
Averaging attention-weighted pixels in \textbf{float32} before\lb
non-linearities and gating against pixel features significantly\lb
improves low-light and motion resilience.
\item \textbf{Honesty Gating Prevents False Readings:} Combining a\lb
prior-free causal HMM tracker with a 5-s median SNR gate\lb
(>= -1.5 dB) separates true cardiac pulses from steady\lb
auto-exposure noise.
\end{bul}

\Subsection{Future Work}
Future work will focus on: (1) multi-dataset training and\lb
cross-dataset evaluation on \textbf{MMPD} , \textbf{PURE} , and \textbf{UBFC-Phys}\lb
to validate generalization across diverse illumination regimes and\lb
mobile sensors; (2) unsupervised test-time personal/domain
\colbreak
% ============================ PAGE 7, LEFT ============================
adaptation using self-supervised spectral sparsity losses (SiNC\lb
[13]) on unlabeled webcam video; and (3) hardware camera\lb
auto-exposure/white-balance lock integration.

\SectionNoNum{REFERENCES}
\vspace{-1.6pt}
{\refsize
\begin{list}{}{\setlength\leftmargin{15.7pt}\setlength\labelwidth{15.0pt}%
  \setlength\labelsep{0.7pt}\setlength\itemsep{1.4pt}\setlength\parsep{0pt}%
  \setlength\topsep{0pt}\setlength\partopsep{0pt}\setlength\itemindent{0pt}%
  \renewcommand\makelabel[1]{\hspace{0.7pt}#1\hfil}}
\item[{[1]}] D. McDuff, S. Gontarek, and R. Picard, "Remote\newline
photoplethysmography for heart rate and blood oxygenation\lb
measurement: A review," \textit{IEEE J. Biomed. Health Inform.} , vol. 19,\lb
no. 6, pp. 1845--1859, 2015.
\item[{[2]}] S. Bobbia, R. Macwan, Y. Benezeth, A. Mansouri, and J. Dubois,\lb
"Unsupervised skin tissue segmentation for remote\newline
photoplethysmography," \textit{Pattern Recognition Letters} , vol. 124, pp.\lb
82--90, 2019.
\item[{[3]}] G. Hsu, A. Ambikapathi, and M. Chen, "Deep learning framework for\lb
heart rate estimation from facial videos," in \textit{Proc. IEEE Int. Conf.\lb
Syst., Man, Cybern. (SMC)} , 2017.
\item[{[4]}] W. Chen and D. McDuff, "DeepPhys: Video-based physiological\lb
measurement using convolutional attention networks," in \textit{Proc. Eur.\lb
Conf. Comput. Vis. (ECCV)} , 2018, pp. 349--365.
\item[{[5]}] Z. Yu, X. Li, and G. Zhao, "Remote photoplethysmograph signal\lb
measurement from facial videos using spatio-temporal networks," in\lb
\textit{Proc. Brit. Mach. Vis. Conf. (BMVC)} , 2019.
\item[{[6]}] W. Wang, A. C. den Brinker, S. Stuijk, and G. de Haan, "Algorithmic\lb
principles of remote PPG," \textit{IEEE Trans. Biomed. Eng.} , vol. 64, no. 7,\lb
pp. 1479--1491, 2017.
\item[{[7]}] X. Liu, J. Fromm, S. Patel, and D. McDuff, "Multi-task temporal shift\lb
attention networks for on-device contactless vitals measurement," in\lb
\textit{Proc. Adv. Neural Inf. Process. Syst. (NeurIPS)} , 2020, pp.\lb
19400--19411.
\item[{[8]}] M. Kumar, A. Veeraraghavan, and A. Sabharwal, "DistancePPG:\lb
Robust non-contact vital signs monitoring using a camera," \textit{Biomedical\lb
Optics Express} , vol. 6, no. 5, pp. 1565--1588, 2015.
\item[{[9]}] G. de Haan and V. Jeanne, "Robust pulse rate from\newline
chrominance-based rPPG," \textit{IEEE Trans. Biomed. Eng.} , vol. 60, no.\lb
10, pp. 2878--2886, 2013.
\item[{[10]}] M. P. Tarvainen, P. O. Ranta-aho, and P. A. Karjalainen, "An\lb
advanced detrending method with application to HRV analysis," \textit{IEEE\lb
Trans. Biomed. Eng.} , vol. 49, no. 2, pp. 172--175, 2002.
\item[{[11]}] X. Niu, S. Shan, H. Han, and X. Chen, "RhythmNet: End-to-end heart\lb
rate estimation from face via spatial-temporal representation," \textit{IEEE\lb
Trans. Image Process.} , vol. 29, pp. 2409--2423, 2020.
\item[{[12]}] Z. Zhang, Z. Pi, and B. Liu, "TROIKA: A general framework for heart\lb
rate monitoring using wrist-type PPG signals during intensive physical\lb
exercise," \textit{IEEE Trans. Biomed. Eng.} , vol. 62, no. 2, pp. 522--531,\lb
2015.
\item[{[13]}] J. Speth, N. Vance, P. Flynn, and A. Czajka, "Non-contrastive\lb
unsupervised learning of physiological signals from video," in \textit{Proc.\lb
IEEE/CVF Conf. Comput. Vis. Pattern Recognit. (CVPR)} , 2023, pp.\lb
14464--14474.
\item[{[14]}] X. Liu, B. Hill, Z. Jiang, S. Patel, and D. McDuff, "rPPG-Toolbox:\lb
Deep remote PPG toolbox," in \textit{Proc. Adv. Neural Inf. Process. Syst.\lb
(NeurIPS)} , 2023.
\item[{[15]}] Z. Sun and X. Li, "Contrast-Phys: Unsupervised video-based remote\lb
physiological measurement via spatiotemporal contrast," in \textit{Proc. Eur.\lb
Conf. Comput. Vis. (ECCV)} , 2022, pp. 492--510.
\item[{[16]}] Z. Yu, Y. Shen, J. Shi, H. Zhao, P. Torr, and G. Zhao, "PhysFormer:\lb
Facial video-based physiological measurement with temporal difference\lb
transformer," in \textit{Proc. IEEE/CVF CVPR} , 2022, pp. 4186--4196.
\end{list}
\par}

\end{document}